\section{Study Limitations}
\label{sec:limitations}

This study is subject to several limitations. The multimodal overlapping cohort contains only $n = 126$ TCGA-KIRC patients with 18 $M_1$ cases, given that genomic, imaging, and clinical alignment requires overlap across TCGA-KIRC inclusion criteria, SEER registry characterization, and CT acquisition \cite{schulz2021multimodal, mota2024multimodal}; fusion-rule gain detection is statistically under-powered and external validation is constrained by the same alignment bottleneck. The retrospective design imposes constraints that would be lifted by prospective validation with standardized CT acquisition and harmonized genomic assays. The 14.29\% $M_1$ prevalence limits positive-class statistical power despite SMOTE correctives; while LOO resampling confirms only modest AUROC variability, the absolute sample size remains the primary limitation. The fusion evaluation is constrained to clinical, genomic, and CT imaging modalities, and other modalities, including histological whole-slide imaging, MRI, or radiopathomics \cite{chi2026multimodal, xiao2026urofusion}, may exhibit different calibration dynamics. External validation on non-TCGA cohorts has not been performed, and cross-institutional CT-acquisition differences in scanner type, slice thickness, and reconstruction kernel may substantially affect radiomic feature distributions. Information-theoretic fusion rules, particularly Dempster--Shafer with contextual discounting and Optimal Transport with Wasserstein-1 minimization, scale poorly to large patient cohorts \cite{huang2024deep, courty2016optimal}; scaling to $n > 1000$ requires approximations that may alter fusion outcomes. Selection bias through TCGA inclusion criteria, batch effects in genomic assays, and operator-dependent radiomic segmentation can each introduce systematic distortions not captured by 5-fold stratified cross-validation alone. Finally, only SMOTE was evaluated as a class-rebalancing strategy; class weighting, focal loss, and threshold tuning may yield different calibration dynamics. The RAG evaluation was conducted on a curated 10-document ccRCC guideline corpus; larger, noisier corpora with cross-institutional guidelines may exhibit different retrieval-calibration dynamics. The Qwen2.5-3B-Instruct model, while producing clinically grounded outputs, is smaller than state-of-the-art medical LLMs; larger models (7B--70B) may exhibit different faithfulness-hallucination tradeoffs.

\section{Future Work}
\label{sec:future_work}

Several concrete research directions are warranted. Larger multicenter ccRCC cohorts with harmonized imaging acquisition would increase alignment sample size and statistical power for fusion-rule gain detection. Prospective collection of clinical, genomic, and CT radiomic data with operationalized calibration protocols would test whether the Calibration--Fusion Hazard persists in real-world deployment. Recent architectures including multimodal Transformers with cross-modal attention, MMIST-ccRCC co-training, and product-of-experts fusion networks deserve explicit calibration-aware re-evaluation. Self-supervised CT encoders such as CT-Foundation, MedSAM, and RadFM may yield richer single-modality representations whose complementarity under calibrated ensembles merits investigation. Bayesian deep ensembles and Dirichlet-based uncertainty quantification methods \cite{shao2023deep, huang2024deep} can be evaluated under the same calibration-aware protocol. Temperature scaling and beta-calibration at the fused-output level represent natural extensions to Platt scaling at the base-model level. SHAP-guided feature attribution should be combined with calibration-aware evaluation to ensure that feature-importance interpretations remain valid under post-processing. Recent frameworks that maintain fusion performance under substantial modality dropout should be re-evaluated under missing-modality conditions, where imputation may introduce further calibration drift. For the RAG dimension, scaling to larger biomedical corpora (PubMed abstracts, clinical trial databases), evaluating larger open-weight medical LLMs (BioMistral-7B, Med-Gemma, Llama-3-Med), and deploying agentic RAG in prospective clinical decision support workflows represent natural extensions. The interaction between retrieval calibration and LLM generation faithfulness under domain shift (e.g., non-KIRC urological cancers) warrants systematic investigation.

\section{Conclusion}
\label{sec:conclusion}

This study performed a rigorously zero-leakage multimodal evaluation of distant metastasis prediction in clear cell renal cell carcinoma across two complementary paradigms: decision-level classifier fusion and LLM-powered retrieval-augmented generation.

\textbf{Stage~I (Decision-Level Fusion):} Three single-modality base models with seven decision-level fusion strategies were benchmarked under three probability calibration regimes on a 126-patient TCGA-KIRC aligned cohort. Single-modality CT radiomics provided the strongest discrimination (AUROC 0.7037). Uncalibrated information-theoretic fusion produced artifactual AUROC gains up to 0.7505 that dissolved under 5-fold Platt cross-calibration (Platt $F_2$-weighted primary endpoint AUROC 0.6631, $p = 0.779$; Platt Dempster--Shafer AUROC 0.6831, $p = 0.442$).

\textbf{Stage~II (LLM-Powered RAG):} A real Qwen2.5-3B-Instruct language model with PubMedBERT dense retrieval was deployed across Naive, Multimodal Hybrid, and Agentic (ReAct tool-calling) RAG architectures on 150 multimodal ccRCC clinical decision queries. Calibrated Multimodal RAG achieved statistically significant superiority over Naive RAG (RAGAS $+0.0481$, $p = 0.001$), confirming genuine retrieval complementarity. Uncalibrated Reciprocal Rank Fusion produced catastrophic probability distortion (ECE = 0.5718), while all architectures achieved zero hallucination rate when grounded on domain-specific clinical guidelines.

The methodological contribution is the formal identification and mechanistic decomposition of the \textbf{Retrieval--Fusion Calibration Hazard}: a unified evaluation artifact affecting both decision-level fusion and RAG systems that operate on heterogeneous, uncalibrated score distributions. The Hazard was decomposed into four independent mechanisms: inter-modality residual error correlation ($\rho \in [0.3842, 0.4619]$), Dempster--Shafer evidence conflict ($\bar{K} = 0.4381$ with 54.8\% high-conflict patients), SMOTE-induced probability stretching (+0.0438 AUROC boost), and retrieval-score mixing artifact (ECE inflation from 0.0633 to 0.5718). The Hazard is not a property of fusion rules or RAG architectures per se, but of fusion composed with uncalibrated, heterogeneous base scores. Per-modality out-of-fold probability calibration before fusion is a mandatory precondition for valid performance reporting. The broader significance extends beyond ccRCC: as biomedical AI increasingly applies deep multimodal architectures and LLM-powered RAG to clinical decision support, the probability-calibration precondition for valid fusion should become a routine, mandatory element of evaluation pipelines. In biomedical AI, a fusion gain that survives Platt cross-calibration is a fusion gain worth reporting.

\section*{Declaration of Funding}
This research received no specific grant from funding agencies in the public, commercial, or not-for-profit sectors.

\section*{Declaration of Generative AI in Scientific Writing}
During the preparation of this work, the authors used generative AI tools solely for linguistic refinement, grammar checking, and structural formatting of the manuscript text. After using these tools, the authors reviewed and edited the content as needed and take full responsibility for the content of the publication. The entire study design, experimental pipeline, code implementation, data curation, statistical evaluations, and scientific interpretations were independently developed and empirically verified by the authors through original computational experiments.

\section*{Data Availability Statement}
The population-scale clinical dataset from the Surveillance, Epidemiology, and End Results (SEER) program contains confidential patient health information and is available upon request following execution of a formal Data Use Agreement with the National Cancer Institute (NCI) SEER program. The genomic RNA-sequencing data and clinical records for the TCGA-KIRC cohort are publicly available from the NCI Genomic Data Commons (GDC; \url{https://portal.gdc.cancer.gov/}). The 3D CT radiomics imaging volumes for the TCIA-KIRC and KiTS23 cohorts are publicly accessible from The Cancer Imaging Archive (TCIA; \url{https://www.cancerimagingarchive.net/}) and the KiTS23 Challenge repository (\url{https://kits-challenge.org/kits23/}). Complete source code, cross-validation evaluation pipelines, tabular metrics, and figure generation scripts for the RenoFusion framework are open-source and publicly available on GitHub at \url{https://github.com/ali-Hamza817/RenoFusion}.

\section*{Declaration of Competing Interest}
The authors declare that they have no known competing financial interests or personal relationships that could have appeared to influence the work reported in this paper.

\section*{CRediT Authorship Contribution Statement}
\textbf{Ali Hamza:} Conceptualization, Methodology, Software, Validation, Formal analysis, Investigation, Data Curation, Writing -- Original Draft, Visualization, Project administration. \\
\textbf{Imran Usman:} Conceptualization, Supervision, Methodology, Resources, Writing -- Review \& Editing, Project administration.
